\documentclass[10pt]{article}
\usepackage[margin=20mm]{geometry}
\usepackage{amsmath,amssymb,fontspec,unicode-math,xurl,hyperref}
\IfFontExistsTF{DejaVu Serif}{\setmainfont{DejaVu Serif}}{\setmainfont{Latin Modern Roman}}
\setmathfont{latinmodern-math.otf}
\setlength{\parindent}{0pt}\setlength{\parskip}{6pt}\sloppy
\begin{document}
\# Scientific verification report: I25-40116\par
\# 1. Identity and exact reviewed version\par
Paper: Position: The Artificial Intelligence and Machine Learning Community Should Adopt a More Transparent and Regulated Peer Review Process. Author: Jing Yang. ICML 2025, PMLR 267:82399-82408. Assigned display number: 0017.\par
The reviewed source is the 10-page final PMLR PDF, SHA-256 `038bff56e4440b9896d64fcc0eb6819cbeb4fd051c48f16f08875b129178723d`, downloaded at 2026-10-24T04:25:43.932578+00:00. Page references below are 1-based PDF pages and match printed internal pages. No appendix is present. [Yang-PMLR (2025)]\par
Task B001 / [internal task] / attempt-001. Batch token: [internal identifier omitted]. Review date: 23 October 2026, America/Chicago.\par
\# 2. Summary\par
Supported outcome: Major. The sixfold active-user comparison and the equation of closed visibility with one-time rebuttal materially compromise two parts of the paper’s central review-model analysis. The U.S. patent-disclosure exception remains a Minor local correction. The causal benefits attributed to openness remain Needs Evidence, because the reported comparisons do not isolate disclosure policy from venue, content, timing, and selection effects.\par
P1 affects the quantitative evidence in Section 4.3, Community Interests Validation; P2 affects the explanatory inference linking public disclosure to iterative author clarification. These are core empirical and procedural dependencies. Major does not require disproving all benefits of open review or reversing the direction of an association. The original full-text, source-policy, visual, and arithmetic checks are retained; acceptance rechecked the relevant source anchors and dependency. The historical website analytics and individual survey responses were not independently reproduced.\par
\# 3. Definitions, data, and scope\par
The paper defines review modes by timing and extent of public disclosure, while its engagement measures include website views, active users, engagement duration, confidence scores, and reply counts. These outcomes are not interchangeable. Website visitors are not automatically a representative sample of researchers; stated confidence does not by itself measure calibration; reply counts do not directly measure review quality.\par
The comparison mixes public API records, voluntary community submissions, and website analytics. Figures 7-9 concern accepted papers. Table 3 identifies differing data sources. The survey in section 5.2 is a voluntary front-page survey about anonymous score sharing. These design details are central to interpreting its claims. [Yang-PMLR (2025), pp.3-8]\par
All ten PDF pages were read. Figures 5-9 were visually checked against relevant narrative statements. Public materials and version histories were inspected, but no paper-specific frozen Google Analytics/Search Console export or complete survey package was located in the inspected links. No author code was run.\par
\# 4. Problem 1: the active-user multiplier does not match Figure 5\par
\#\# Original claim and location\par
On p.5, section 4.3, the text gives ICLR 88,220 active users and describes this as approximately six times the NeurIPS count. Figure 5 on the same page displays active users as orange bars against the left logarithmic count axis. [Yang-PMLR (2025), p.5]\par
\#\# Complete numerical and visual check\par
The exact source page was rendered with pdftoppm. The NeurIPS orange bar is unambiguously above 10\textasciicircum{}4.5, or approximately 31,623, rather than near 14,700. Only this conservative graphical bound is used:\par
{\ttfamily\footnotesize NeurIPS count implied by sixfold statement = 88,220 / 6 = 14,703.3333.\\
Conservative lower bound from plot = 10\textasciicircum{}4.5 = 31,622.7766.\\
Therefore ICLR / NeurIPS < 88,220 / 31,622.7766 = 2.7897614.}\par
Thus the sixfold comparison and the plotted count cannot both describe the same active-user metric, reporting window, and filter. Reading the wording as sevenfold rather than sixfold does not resolve the disagreement. The figure crop is preserved as sources/0017/figure5\_crop.png, derived from the hashed original. [Independent arithmetic; independent\_results.json]\par
\#\# Counterevidence, impact, and repair\par
Possible explanations include different reporting windows, filters, or a mistaken comparator. None is specified in the compared text and caption. The check does not establish which representation contains the correct raw count.\par
Judgment: Major. Section 4.3 presents this comparison as quantitative validation of community interest across review modes. Under the common measure and population implied by the text and caption, a claimed factor of six is incompatible with the conservative graphical upper bound below 2.79. This materially overstates a core empirical comparison, even though the displayed direction of the association survives. The bound establishes inconsistency, not which raw representation is correct. A precise numerical repair requires both counts for the same population, window, and filters; no exact replacement ratio is inferred from pixel heights. Repair verification remains pending.\par
\# 5. Problem 2: private review is conflated with one-time author response\par
\#\# Original claim and suspect step\par
At the opening of p.7, section 4.3, the discussion contrasts extended open-review dialogue with a categorical characterization of fully closed models as allowing a single rebuttal phase. The suspect implication is that non-public reviews necessarily prevent iterative author-reviewer communication. [Yang-PMLR (2025), p.7]\par
\#\# Complete check against a historical counterexample\par
The paper classifies ICML in its closed category for the historical period ending in 2024. The official ICML 2024 Author Instructions, under Reviewing and Author Response, specify private reviewing, an initial rebuttal, a subsequent discussion period involving authors, reviewers and area chairs, and permission for multiple responses per submission. This is an example within the paper’s own historical scope, rather than a policy introduced only after publication. [ICML-2024, Reviewing and Author Response]\par
The relevant variables are distinct: public visibility, number of allowed responses, duration of discussion, and eligibility to participate. A private process can allow multiple exchanges, while excluding the broader public.\par
\#\# Counterevidence, impact, and precise repair\par
Open public commenting can still expand participation beyond the author/reviewer group. Nothing here establishes equal interaction volumes or quality across review modes. A charitable reading is that the text intended a particular closed venue, but it did not specify that restriction.\par
Judgment: Major. The paper defines review modes by public disclosure on p.1 and places ICML in the fully closed category in Figure 5. The within-scope ICML 2024 counterexample therefore invalidates the categorical premise used in Section 4.3 to explain iterative clarification as an advantage of open over closed review. Public visibility and permission for repeated author responses are distinct policy variables. This is a failed core explanatory inference, not an assessment that public participation has no value. Repair by identifying venue/year, response rounds, discussion duration, and eligible participants separately from disclosure. The policy evidence verifies this distinction; it does not estimate the contribution of each variable to the observed reply differences.\par
\# 6. Problem 3: the U.S. patent statement is too absolute\par
\#\# Original claim and location\par
Section 6, p.8, Disclosure Policy, describes the U.S. first-to-file system as requiring no public disclosure before filing. The adjacent statement that an open submission can jeopardize patent protection is conditional and broader. [Yang-PMLR (2025), p.8]\par
\#\# Source check and corrected scope\par
USPTO MPEP section 2153.01, revision R-01.2024, explains 35 U.S.C. 102(b)(1)(A): qualifying inventor-originated disclosures made no more than one year before the effective filing date are excepted from the specified prior-art rule. This pre-existing exception disproves an unconditional U.S. prohibition on all prefiling disclosure. [USPTO-2153.01 (2024)]\par
For a concrete logical instance, assume an otherwise qualifying inventor-originated disclosure occurs six months before the effective filing date and satisfies the statutory provenance and subject-matter conditions. The paper’s absolute statement would exclude it solely because disclosure preceded filing; the cited exception does not treat that qualifying disclosure as prior art under 102(a)(1). This does not establish patentability under all other requirements.\par
Judgment: Minor. The defect is local to the legal premise used in the confidentiality discussion. A defensible repair acknowledges the defined U.S. exception and conditions the risk on jurisdiction, timing, disclosure origin, and subject matter. The general concern about disclosure risk remains. No particular invention, application, or filing strategy was evaluated.\par
\# 7. Problem 4: causal effects of openness remain unidentified\par
\#\# Claim and location\par
Sections 4.3 and 5, pp.5-8, attribute engagement and oversight advantages to review openness. Website traffic and accepted-paper discussion distributions are offered as supporting observations. [Yang-PMLR (2025)]\par
\#\# Check of the inferential step\par
The comparisons vary venue and review mode together. Venue identity also carries topic, audience, calendar, available content, search visibility, and review procedures. The site’s ability to show live information differs by mode. Restricting discussion data to accepted papers introduces venue-specific selection. Similar submission counts do not remove these differences.\par
A minimal identification argument makes the limit explicit. Let m\_c be an observed mean engagement value for venue c and O\_c its disclosure indicator. In the decomposition m\_c = beta*O\_c + gamma\_c, gamma\_c represents unrestricted venue effects. For any beta, choosing gamma\_c = m\_c - beta*O\_c reproduces every observed mean. The reported cross-venue means therefore do not identify beta without additional assumptions or design information. This is an analytic identification check, not a fitted model or evidence of a zero effect.\par
\#\# Alternative readings and supported impact\par
The paper is a position paper and can legitimately use descriptive patterns to motivate policy. Some explanations are tentative, and p.9 acknowledges voluntary-selection bias. Its traffic data can support interest in the particular information service. Those points limit the concern.\par
Judgment: Needs Evidence. The direction or size of a causal effect, the quality of discussion, the reduction of overconfidence, and population-wide policy preferences are unresolved. They are not elevated into supported Major damage. A verified wording repair would describe associations and hypotheses. A causal repair requires a defensible design or assumptions and relevant data, not merely a new label.\par
\# 8. Additional checked claims and scope clarifications\par
The basic CVPR 2025 author-reviewer obligation is documented. Its official policy also limits assignments to qualified reviewers, varies load with seniority, and lists exemptions. Therefore the obligation itself is not fabricated, while the paper’s proposed quality consequences require evidence that considers these safeguards. The text already calls some complaints anecdotal. [CVPR-2025; Yang-PMLR (2025), p.7]\par
The 57\% anonymous-score-sharing result is a statement about voluntary respondents, not a representative estimate of support for full, real-time public reviews. The observed age distribution also does not establish every visitor’s occupation or career status. These are interpretation limits, not demonstrations that the raw reported values are false.\par
The lower-confidence explanation on p.6 is explicitly tentative. Confidence calibration would require comparison with an appropriate correctness or quality outcome. Visual inspection found that the stated reply maxima and increasing ICLR medians are consistent with Figures 8-9; those displays do not independently establish improved rigor.\par
The four Table 2 percentages sum to 100.8\%. This cannot be explained by ordinary one-decimal rounding if they are mutually exclusive exhaustive shares, but overlapping users or differing attribution definitions might explain it. The denominator is unresolved, so this was not promoted to a confirmed arithmetic defect.\par
Figure 3 explicitly excludes 2025. ICML’s 2025 FAQ describes later publication of accepted-paper reviews, with opt-in publication for rejected papers. That policy change is recorded as time-sensitive context rather than retroactively invalidating the figure’s historical classification. [ICML-2025-FAQ]\par
\# 9. Later versions, public materials, and overall dependence\par
The arXiv record lists versions v1 (2 February 2025), v2 (24 May 2025), and v3 (15 October 2025). A support agent independently compared extracted text from v2 and v3 with the PMLR file, joining line-break hyphenation, removing the repeated running header, and collapsing whitespace. The differences were the title prefix, arXiv stamp, and added section numbering. No substantive textual repair of the inspected claims was found. Images were not compared across versions, so complete visual identity is not asserted. The downloaded v3 has SHA-256 `f418e3cc8b1213622b7ded7a4fc9f99d9badaeba440779ac2fc6e2a34164c4b0`. [Yang-arXiv-v3]\par
Public Paper Copilot repositories contain conference and review records, but the documentation inspected did not identify the historical website-analytics exports, demographic missingness/denominators, or complete survey records required to reproduce this paper’s traffic and attitude claims. Current repository contents and later website usage totals are not substitutes for that frozen historical snapshot. [Paperlists; OpenReview-archive; ContactUs]\par
No matching author erratum was identified in the inspected version record, site documentation, or scoped project-issue search. OpenReview discussion was inaccessible. The separate later Paper Copilot system paper was used only to locate released resources, not reviewed as an additional assigned paper.\par
The contribution-dependency disposition is explicit: P1 is Major because an internally incompatible magnitude supports the central empirical comparison of review modes; P2 is Major because a false categorical procedural premise supports the interpretation of iterative discussion benefits. P3 remains Minor because its qualified patent-law correction concerns an alternative-view example and does not supply the empirical or procedural comparison. P4 remains Needs Evidence: the data do not identify a causal policy effect, and no effect reversal is claimed. The normative recommendations, public-participation distinction, and descriptive observations can survive these corrections. Their survival does not lower the severity of the failed core steps. Human policy and scientific decisions remain pending.\par
\# 10. Conclusion table\par
Item | Judgment | Repair verification | Adoption | Remaining work\par
Active-user multiplier | Major; core magnitude inconsistent | Bound checked; exact repair pending | Not requested | Original counts, common window/filter\par
Closed review and rebuttals | Major; core premise false | Policy-based correction checked | Not requested | Venue-specific comparative analysis\par
U.S. prefiling disclosure | Minor; statutory exception omitted | Limited legal correction sourced | Not requested | No specific patent case assessed\par
Openness causes better outcomes | Needs Evidence | Descriptive wording defensible | Not requested | Causal design and outcome measures\par
Survey, confidence, historical rules | Scope clarification | Definitions and primary policies checked | Not applicable | Raw exports and population inference\par
\# 11. Consulted references\par
[Yang-PMLR (2025)] Jing Yang. PMLR 267:82399-82408. https://proceedings.mlr.press/v267/yang25ba.html . Matching immutable PDF: https://raw.githubusercontent.com/mlresearch/v267/a3e4112da70e3f0623daeb3eeea06e2b34ed5e2f/assets/yang25ba/yang25ba.pdf .\par
[ICML-2024] Official Author Instructions, Reviewing and Author Response. https://icml.cc/Conferences/2024/AuthorInstructions .\par
[USPTO-2153.01 (2024)] Manual of Patent Examining Procedure, section 2153.01, R-01.2024; 35 U.S.C. 102(b)(1)(A). https://www.uspto.gov/web/offices/pac/mpep/s2153.html .\par
[CVPR-2025] Official CVPR 2025 Changes, author-reviewer obligation, qualification screening and exceptions. https://cvpr.thecvf.com/Conferences/2025/CVPRChanges .\par
[ICML-2025-FAQ] Official Peer Review FAQ, post-decision questions. https://icml.cc/Conferences/2025/PeerReviewFAQ .\par
[Yang-arXiv-v3] arXiv:2502.00874v3, 15 October 2025. https://arxiv.org/abs/2502.00874v3 and https://arxiv.org/html/2502.00874v3 .\par
[Paperlists] https://github.com/papercopilot/paperlists , inspected documentation at HEAD `9fcfd1234879a8e5630b5116dde411ec143dd763`. [OpenReview-archive] https://github.com/papercopilot/openreview , inspected documentation at HEAD `3bb756daff86e9b084ada9f1c9882da108a1be06`. [ContactUs] https://github.com/papercopilot/ContactUs , inspected documentation at HEAD `f15e2695d29a697750bed96472e8b9185b4cb986`.\par
\# Appendix A. Commands, finite results, and provenance\par
The original download and text extraction are recorded in sources/0017/source\_download.json. Supporting official HTML snapshots and arXiv v3 are hashed in sources/0017/supporting\_downloads.json. The local independent script sources/0017/independent\_checks.py executed successfully and wrote verification/0017/independent\_results.json.\par
Actual independent calculation expressions:\par
{\ttfamily\footnotesize 88220 / 6\\
10 ** 4.5\\
88220 / (10 ** 4.5)\\
sum([59.9, 23.9, 9.4, 7.6])}\par
Results, in order: 14703.333333333334; 31622.776601683792; 2.789761351800544; 100.8. The graphical inequality depends on visual inspection of the labeled logarithmic axis, not on assumed raw analytics. No model training, author-code execution, or empirical reproduction occurred.\par
The original page-5 rendering command used pdftoppm -f 5 -l 5 -singlefile -r 160 -png with the assigned source PDF. The preserved crop is evidence for the visual comparison. A support agent’s direct Lumen PDF reader failed with a toLowerCase property error; pdftotext and pdftoppm processed the exact source successfully.\par
\# Appendix B. Roles, access limits, and rendering\par
The parent assistant performed synthesis, source-policy verification, visual checking, and independent arithmetic. AI support agents separately checked methodology, figures/data scope, and version history. The selected model label in the task is 6 Pro; the assistant environment identifies itself as GPT-6 Astra Pro. The controller’s visible UI model field was not inspected. No human-verification declaration or author adoption is asserted.\par
OpenReview returned a browser-verification wall. No private analytics account was accessed. The exact historical analytics and survey exports were not located in the inspected public materials, and the underlying published data-collection pipeline was not run. There was no separate supplement in the downloaded PDF or listed PMLR related materials; inaccessible OpenReview attachments remain unverified.\par
The original and acceptance-proposal artifacts are preserved with predecessor hash maps. The acceptance proposal supplied a separately rendered preview and requested controller compilation. This version uses a PDF compiled from the accompanying TeX with the existing desktop Tectonic compiler; exact build hashes and visual checks are retained in verification. No author code was executed and no software was installed. Artifact acceptance and version promotion do not constitute human scientific adjudication; scientific decisions remain pending.\par
\end{document}
